% =============================================================================
%  Capítulo 3 — Desarrollo de Software y Diseño
%
%  Este capítulo está casi vacío en tu PDF actual. Aquí tienes el esqueleto
%  completo según la rúbrica del enunciado.
%
%  RECUERDA: el "Capítulo 4 Desarrollo" vacío de tu PDF actual NO existe en
%  esta plantilla. La rúbrica solo prevé un capítulo de "Desarrollo de software".
% =============================================================================
\chapter{Desarrollo de Software y Diseño}
\label{cap:desarrollo}

% -----------------------------------------------------------------------------
\section{Análisis de Requisitos}
\label{sec:requisitos}

Los siguientes requisitos se han definido de manera específica para permitir
la ejecución de tests automatizados que demuestren su cumplimiento.

\subsection{Requisitos funcionales}
\label{sec:rf}

\todo{Lista numerada con referencias al §2. Ejemplos:
\begin{itemize}
  \item \textbf{RF-01:} El sistema debe implementar el proceso de difusión VE
        descrito en §\ref{sec:ve}.
  \item \textbf{RF-02:} El sistema debe implementar el proceso de difusión VP
        descrito en §\ref{sec:vp}.
  \item \textbf{RF-03:} El sistema debe permitir seleccionar entre los procesos
        VE y VP mediante un parámetro de configuración.
  \item \textbf{RF-04:} El sistema debe implementar el sampler Euler--Maruyama.
  \item \textbf{RF-05:} El sistema debe implementar el sampler
        Predictor--Corrector con SNR objetivo configurable.
  \item \textbf{RF-06:} El sistema debe implementar el sampler Probability
        Flow ODE.
  \item \textbf{RF-07:} El sistema debe implementar las schedules lineal y
        coseno.
  \item \textbf{RF-08:} El sistema debe calcular las métricas BPD, FID e IS
        sobre un conjunto de muestras.
  \item \textbf{RF-09:} El sistema debe permitir la generación condicionada por
        clase mediante CFG.
  \item \textbf{RF-10:} El sistema debe permitir la imputación de imágenes
        parcialmente observadas mediante una máscara binaria.
  \item \textbf{RF-11:} El sistema debe soportar imágenes en blanco y negro
        (1 canal) y en color (3 canales).
\end{itemize}}

\subsection{Requisitos no funcionales}
\label{sec:rnf}

\paragraph{Entorno y herramientas.}
\todo{Versiones exactas. Mirar \texttt{.python-version} y \texttt{uv.lock}:
\begin{itemize}
  \item Python: \todo{X.Y}
  \item PyTorch: \todo{X.Y}
  \item Gestor de dependencias: \texttt{uv}
  \item Otras: NumPy, scipy, matplotlib, torchvision, scikit-image, etc.
\end{itemize}}

\paragraph{Hardware.}
\todo{GPU usada (modelo, VRAM), CPU, RAM. Indicar mínimos para
inferencia y para entrenamiento por separado.}

\paragraph{Rendimiento.}
\todo{Especificaciones cuantitativas verificables, p.ej.:
"Una imagen 28$\times$28 debe generarse en menos de 2~s con el sampler
Euler--Maruyama y $N=500$ pasos en una GPU NVIDIA RTX 3060". Cada uno
debe corresponderse con una entrada de la tabla de Resultados §\ref{cap:resultados}.}

\paragraph{Compatibilidad.}
\todo{Sistemas operativos soportados, formatos de checkpoint
(\texttt{.pth}), interoperabilidad con otros frameworks.}

\paragraph{Estilo y documentación.}
\todo{PEP-8 verificado con \texttt{ruff}/\texttt{black}, type hints,
docstrings estilo NumPy, generación automática con Sphinx/mkdocs.}

\paragraph{Metodología de desarrollo.}
\todo{Git/GitHub, ramas \emph{feature}, pull requests, revisión cruzada,
CI con GitHub Actions.}

% -----------------------------------------------------------------------------
\section{Casos de Uso}
\label{sec:casos-uso}

\todo{Describir los casos de uso principales, cada uno enlazado a un
notebook real del proyecto. Ejemplos:
\begin{itemize}
  \item \textbf{CU-01 — Generación incondicional}: usuario entrena un modelo
        VP sobre MNIST y muestrea N imágenes nuevas.
        Notebook: \texttt{project\_AAIII\_teamCode.ipynb}.
  \item \textbf{CU-02 — Generación condicionada por clase}: usuario carga
        modelo CFG entrenado sobre SVHN y genera imágenes de la clase 7.
        Notebook: \texttt{cfg\_diffusion\_models\_SVHN.ipynb}.
  \item \textbf{CU-03 — Imputación}: usuario carga imagen parcialmente
        observada y completa los píxeles faltantes.
  \item \textbf{CU-04 — Comparativa de samplers}: usuario evalúa FID de
        EM, PC y PF-ODE sobre el mismo modelo.
  \item \textbf{CU-05 — Cálculo de BPD}: usuario evalúa BPD sobre un
        conjunto de test.
\end{itemize}}

% -----------------------------------------------------------------------------
\section{Arquitectura y Diseño}
\label{sec:arquitectura}

\subsection{Estructura del paquete}
\label{sec:estructura-paquete}

El sistema se organiza como un paquete Python instalable con la siguiente
estructura de módulos:

\begin{lstlisting}[language=, caption={Estructura del paquete \texttt{diffusion\_lib}}]
diffusion_lib/
├── __init__.py
├── model.py                # U-Net (score network)
├── processes/              # Procesos de difusión
│   ├── base.py             # BaseProcess (ABC)
│   ├── ve.py               # Variance Exploding
│   └── vp.py               # Variance Preserving
├── samplers/               # Integradores de la SDE inversa
│   ├── base.py             # BaseSampler (ABC)
│   ├── euler_maruyama.py
│   ├── predictor_corrector.py
│   ├── probability_flow_ode.py
│   └── imputation.py
├── schedules/              # Noise schedules
│   ├── base.py             # BaseSchedule (ABC)
│   ├── linear.py
│   └── cosine.py
└── metrics/                # Métricas de calidad
    ├── bpd.py              # Bits per dimension
    └── fid_is.py           # FID e IS sobre InceptionV3
\end{lstlisting}

\subsection{Diagrama de clases}
\label{sec:diagrama-clases}

\todo{Insertar UML simple. Sugerencia: usar PlantUML o draw.io y exportar
a PDF. Mostrar:
\begin{itemize}
  \item \texttt{BaseProcess} $\leftarrow$ \texttt{VE}, \texttt{VP}
  \item \texttt{BaseSampler} $\leftarrow$ \texttt{EulerMaruyama},
        \texttt{PredictorCorrector}, \texttt{ProbabilityFlowODE}, \texttt{Imputation}
  \item \texttt{BaseSchedule} $\leftarrow$ \texttt{Linear}, \texttt{Cosine}
  \item Relaciones de composición (Sampler usa Process, Process usa Schedule).
\end{itemize}}

% \begin{figure}[H]
%   \centering
%   \includegraphics[width=0.9\textwidth]{diagrama_clases.pdf}
%   \caption{Diagrama de clases del paquete \texttt{diffusion\_lib}.}
%   \label{fig:diagrama-clases}
% \end{figure}

\subsection{Diagrama de secuencia (entrenamiento y muestreo)}
\todo{Opcional pero recomendado. Mostrar el flujo:
\begin{enumerate}
  \item Entrenamiento: dataset $\to$ Process.add\_noise $\to$ U-Net.forward $\to$ loss $\to$ optimizer.
  \item Muestreo: prior $\to$ Sampler.step (loop) $\to$ Process.reverse\_drift $\to$ U-Net $\to$ imagen.
\end{enumerate}}

% -----------------------------------------------------------------------------
\section{Validación y Calidad del Software}
\label{sec:calidad}

\subsection{Tests automatizados}
\todo{Estrategia: \texttt{pytest} con tests unitarios por módulo
(\texttt{test\_processes.py}, \texttt{test\_samplers.py}, etc.) y al menos
un test de integración. Indicar cobertura.}

\subsection{Integración continua}
\todo{GitHub Actions: workflow con steps de \texttt{ruff check},
\texttt{pytest}, build de docs.}

\subsection{Estilo y documentación}
\todo{PEP-8, \texttt{black}, \texttt{ruff}, docstrings NumPy, Sphinx para
auto-documentación de la API.}

\subsection{Licencia}
\todo{Decidir y justificar: MIT (permisiva) / Apache-2.0 (incluye
patentes) / GPL (copyleft). Para un proyecto académico/open-source
suele recomendarse MIT.}

\subsection{Implicaciones éticas}
\todo{Riesgos de modelos generativos: deepfakes, sesgos heredados de los
datasets (MNIST/SVHN/gatos no son problemáticos, pero conviene
mencionar el riesgo en datasets faciales o con sesgo demográfico),
coste energético del entrenamiento.}

% -----------------------------------------------------------------------------
\section{Planificación del Proyecto}
\label{sec:planificacion}

\todo{Diagrama de Gantt. Fases sugeridas:
\begin{enumerate}
  \item Investigación bibliográfica y revisión teórica (semanas 1--3).
  \item Implementación de procesos VE/VP y entrenamiento básico (semanas 3--5).
  \item Implementación de samplers (semanas 4--6).
  \item Implementación de schedules y métricas (semanas 6--7).
  \item Generación condicionada e imputación (semanas 7--9).
  \item Generación de resultados y comparativas (semanas 9--11).
  \item Redacción del informe (semanas 8--12, paralelo).
\end{enumerate}
Recomendado: usar la plantilla de TikZ \texttt{pgfgantt} o exportar
desde una hoja de cálculo.}

% \begin{figure}[H]
%   \centering
%   \includegraphics[width=\textwidth]{gantt.pdf}
%   \caption{Diagrama de Gantt del proyecto.}
%   \label{fig:gantt}
% \end{figure}
